\chapter{Disks, Files, and Pages}
\label{ch:pages}

Every design decision in a storage engine starts from one fact: memory is fast
and volatile, and storage is slow and durable. A database must keep its data on
the slow, durable device and still answer queries at the speed of the fast,
volatile one. This chapter is about the unit that moves between the two: the
\term{page}.

\section{The storage hierarchy}

\begin{table}[ht]
\centering
\begin{tabular}{@{}lll@{}}
\toprule
Device & Typical access latency & Survives power loss \\
\midrule
CPU cache & 1--10 ns & no \\
Main memory (DRAM) & about 100 ns & no \\
NVMe SSD, random 4 KiB read & 10--100 $\mu$s & yes \\
Spinning disk, random read & 5--10 ms & yes \\
\bottomrule
\end{tabular}
\caption{Orders of magnitude, not benchmarks. The gap between memory and storage is a factor of a hundred to a hundred thousand.}
\label{tab:hierarchy}
\end{table}

Two properties of storage shape everything above it.

\begin{description}[leftmargin=1.2em]
\item[Block granularity.] A device reads and writes whole blocks (typically 512 bytes or 4 KiB). Reading one byte costs a whole block. So the engine should arrange data so that one block read brings many useful bytes.
\item[Sequential beats random.] Even on SSDs, sequential access is cheaper than random access, and on spinning disks it is dramatically cheaper. An engine that writes sequentially (a log) and reads mostly through an index (a tree) uses the device the way it wants to be used.
\end{description}

\section{Pages}

A \term{page} is the engine's unit of I/O: a fixed-size block of the database
file, usually 4, 8, or 16 KiB, aligned to the device's block size. The engine
never reads half a page and never writes half a page. Fixing the size has three
payoffs:

\begin{itemize}
\item \textbf{Addressing is arithmetic.} Page $i$ lives at byte offset $i \times \text{PAGE\_SIZE}$ in the file. A page id is all the engine needs to find it.
\item \textbf{Memory is interchangeable.} Every buffer-pool frame holds exactly one page, so any frame can hold any page (\cref{ch:buffer}).
\item \textbf{Failure is bounded.} A write either lands a whole page, or, if power fails mid-write, produces a \term{torn page}: part old, part new. Because the size is fixed and known, a checksum over the page detects the tear.
\end{itemize}

\begin{keyidea}
A page id is a stable name. Pointers inside the database (a B+ tree child, a
heap record's location) are page ids, never memory addresses, so they remain
valid across restarts and across evictions.
\end{keyidea}

\section{Slotted pages}

Fixed-size records are easy to place: record $k$ starts at byte $k \times
\text{size}$. Real data is variable-sized: an IRI might be 20 bytes or 200.
The standard answer is the \term{slotted page}.

\begin{figure}[ht]
\centering
\begin{tikzpicture}[x=1cm,y=1cm,font=\scriptsize]
\fill[boxbg] (0,0) rectangle (1.4,1.2); \node at (0.7,0.6) {header};
\fill[ideabg] (1.4,0) rectangle (3.4,1.2); \node[align=center] at (2.4,0.6) {slot directory\\grows $\rightarrow$};
\node at (6.4,0.6) {free space};
\fill[statusbg] (9.2,0) rectangle (11.0,1.2); \node at (10.1,0.6) {record 2};
\fill[statusbg] (11.0,0) rectangle (12.2,1.2); \node at (11.6,0.6) {rec 1};
\fill[statusbg] (12.2,0) rectangle (14,1.2); \node at (13.1,0.6) {record 0};
\draw (9.2,0) -- (9.2,1.2); \draw (11.0,0) -- (11.0,1.2); \draw (12.2,0) -- (12.2,1.2);
\draw (1.4,0) -- (1.4,1.2); \draw (3.4,0) -- (3.4,1.2);
\draw (0,0) rectangle (14,1.2);
\node at (10.1,1.45) {$\leftarrow$ records grow};
\foreach \sx/\tx in {1.8/13.1,2.4/11.6,3.0/10.1}{\draw[-{Stealth[length=1.6mm]},gray] (\sx,0) .. controls (\sx,-0.7) and (\tx,-0.7) .. (\tx,0);}
\end{tikzpicture}
\caption{A slotted page. The slot directory grows forward from the header, records grow backward from the end, and free space sits between them.}
\label{fig:slotted}
\end{figure}

The header records how many slots exist and where free space begins. Each slot
holds a record's offset and length. A record is addressed by a \term{record id}
(RID): \emph{(page id, slot number)}. Two properties make this work:

\begin{itemize}
\item \textbf{Records can move without changing their RID.} To compact the page, slide records together and update the slot offsets. The slot number, and therefore the RID, is unchanged.
\item \textbf{Deletion is cheap.} Mark the slot free; a later insert may reuse it. The space is reclaimed when the page is compacted, or when a new record fits in the hole.
\end{itemize}

\begin{ontodb}[: pages and the term heap]
Every OntoDB page is 4096 bytes (\code{PAGE\_SIZE}). The first 8 bytes of
\emph{every} page, from the very first version, are reserved for the
\term{pageLSN}, the log sequence number of the last update applied to that
page; the payload starts at byte 8. Reserving the field before any logging
existed means the file format did not have to change when the write-ahead log
arrived (\cref{ch:wal}).

The spellings of IRIs and literals live in a \term{term heap} made of slotted
pages. The B+ tree leaves never contain strings: they contain 24-byte keys of
three ids. That separation keeps the tree's fanout high, and it is why a long
IRI costs nothing in the index (\cref{ch:rdf-indexing}). A spelling is never
deleted when a triple using it is deleted: the dictionary keeps the id, so a
log record written earlier can still name it.
\end{ontodb}

\section{Talking to the file system}

An engine reads and writes its file through the operating system, and two
details of that conversation decide whether data survives a crash.

\paragraph{The page cache.} A \code{write} call returns as soon as the bytes are
in the operating system's page cache, in memory. They reach the device later.
A crash after \code{write} returns can still lose them.

\paragraph{\code{fsync}.} Calling \code{fsync} on a file asks the operating
system to push that file's dirty data to the device and wait until the device
reports it durable. It is the only promise of durability the engine gets, and
it is expensive: a few milliseconds on a disk, tens to hundreds of microseconds
on an SSD. Two consequences follow.

\begin{enumerate}
\item \textbf{Batch the expensive call.} Engines amortize \code{fsync} across many transactions; the log manager's \term{group commit} (\cref{ch:wal}) exists for this.
\item \textbf{Never assume a failed \code{fsync} can be retried.} On several operating systems, after \code{fsync} reports an error the kernel may already have discarded the dirty pages and cleared the error, so a second \code{fsync} that ``succeeds'' has lost data. PostgreSQL learned this in 2018 and now treats a failed \code{fsync} as fatal \cite{fsyncgate,pillai2014}.
\end{enumerate}

\paragraph{Atomic replacement.} Renaming a file over another is atomic on POSIX
file systems: a reader sees the old file or the new one, never a mixture. The
idiom ``write to a temporary file, \code{fsync} it, \code{rename} it over the
target, \code{fsync} the directory'' replaces a whole file atomically. OntoDB
writes its snapshots this way (\cref{ch:durability}).

\begin{decision}{4 KiB pages}
\dpart{Choice} \code{PAGE\_SIZE} is 4096 bytes, fixed for the whole system.
\dpart{Alternatives} 8 KiB (PostgreSQL), 16 KiB (InnoDB), or a size chosen per
database at creation.
\dpart{Why} 4 KiB matches the operating system's page and a common device
block, so one page is one block. It is also small enough that tests with a
pool of three or four frames, and trees of a few hundred keys, exercise
eviction, splits, and merges quickly.
\dpart{Cost} Lower fanout than a 16 KiB page, so a very large index may need
one more level; and fixed per-page overhead (the 8-byte pageLSN, the header) is
a larger share of each page.
\end{decision}

\begin{decision}{Page 0 is a real page}
\dpart{Choice} Page ids are \code{int32}, and page 0 is a legal page.
``No page'' is \code{INVALID\_PAGE\_ID}, which is $-1$. The catalog (plan 3.3) will
initialize its page explicitly, so a zero-filled file cannot be read as ``the root is page
0''.
\dpart{Alternatives} Reserve page 0 as a null page and use 0 as the sentinel.
\dpart{Why} Zero is the value memory and files have by accident. A sentinel
should be a value that never appears unless someone wrote it on purpose; then a
missing initialization is a visible error, not a tree rooted at page 0.
\dpart{Cost} Every reader of a stored page id must check for $-1$, and the
largest file is $2^{31}$ pages, 8 TiB at 4 KiB.
\end{decision}

\begin{decision}{The file's length records allocation}
\dpart{Choice} \code{AllocatePage} grows the database file by one zero page
(\code{ftruncate}, which usually writes no data blocks). On open, the number of
allocated pages is the file size divided by \code{PAGE\_SIZE}, rounded up. There
is no header page of the disk manager's own.
\dpart{Alternatives} A header page (a superblock) holding the next page id and
a free list, as SQLite's first page and InnoDB's space header do.
\dpart{Why} Page 0 already belongs to the catalog. The file system updates a
file's length as one fact, so there is no second record of allocation that a
crash could leave disagreeing with the file.
\dpart{Cost} The new length is durable only after an \code{fsync}, so after a
crash a page that a log record names may lie past the end of the file. The disk
manager therefore reads any page past the end as zeros and grows the file when
one is written there, which lets redo apply such a record with no special case.
\end{decision}

\begin{decision}{The file only grows}
\dpart{Choice} \code{DeallocatePage} does not return the id for reuse. A freed
page stays in the file, unused.
\dpart{Alternatives} An in-memory free list, or a durable free-space map such
as PostgreSQL's.
\dpart{Why} An id that is reused while anything still names it (a log record
awaiting redo, a stale parent pointer after a crash) turns an old reference
into a write to someone else's page. It is the same argument as for term ids:
an id's meaning should not change.
\dpart{Cost} Space freed by merges and deletes is not reclaimed. Reclaiming it
safely needs a logged free-space map, which is later work.
\end{decision}

\begin{decision}{Through the page cache; \code{fsync} at shutdown}
\dpart{Choice} Pages are read and written with \code{pread} and \code{pwrite}
through the operating system's page cache. \code{WritePage} does not
\code{fsync}; \code{ShutDown} does, once.
\dpart{Alternatives} \code{O\_DIRECT}, bypassing the page cache as InnoDB can;
or \code{fsync} after every page write.
\dpart{Why} Durability is the log's job: a committed change survives because
its log record is forced (\cref{ch:wal}), not because its page is. A page needs
to be on disk only at a checkpoint or a clean shutdown. \code{pread} and
\code{pwrite} carry their own offset, so threads share one file descriptor
without a lock around a seek.
\dpart{Cost} A cached page exists twice, in a pool frame and in the kernel's
cache. Until page-level logging exists (plan 8.3), only a clean shutdown makes
the page store durable.
\end{decision}

\begin{decision}{Spellings are never deleted}
\dpart{Choice} Deleting the last triple that uses a term keeps the term and its
id in the dictionary.
\dpart{Alternatives} Reference-count terms and reclaim an id when its count
reaches zero; or garbage-collect unused terms at checkpoint.
\dpart{Why} Log records name terms by id. If id 18 could be reclaimed and
reused, an old record that says ``insert (18, 9, 4)'' would silently mean a
different triple on replay. Never deleting makes an id's meaning permanent.
\dpart{Cost} The dictionary only grows. A workload that churns through unique
literals keeps every one of them.
\end{decision}

\begin{status}
The page frame and the pageLSN layout exist, and the disk manager runs (plan
1.1): it allocates pages by growing the file, reads and writes whole pages, and
fsyncs on shutdown. The term heap is plan 3.1.
\end{status}

\begin{checkyourself}
\item Page 37 of a 4 KiB-page file: at what byte offset does it start, and why does the engine never need a table to find it?
\item In a slotted page, a record is moved during compaction. Why does no index entry that points to it need to change?
\item A \code{write} returned successfully and then the machine lost power. Was the data durable? What call would have made it so?
\item \code{fsync} returns an error. Why is calling it again and continuing on success wrong?
\item Why is $-1$ a better ``no page'' sentinel than 0?
\item After a crash, a log record names page 812, but the file holds only 800 pages. How can OntoDB's allocation scheme produce this, and who must repair it?
\end{checkyourself}

\section*{Further reading}
Petrov \cite{petrov}, chapters 2 and 3, on file formats and slotted pages.
Pillai et al.\ \cite{pillai2014} on what file systems actually guarantee.
